% ECONOMICS_PROBLEM: {"id": "E10-68", "title": "Sufficient heterogeneity in macro models", "jel_code": "E10", "source_id": "OP-0234"}
% FIRST_POSTED: 2026-10-10
% ATTEMPT_STATUS: unresolved
% ENTRY_KIND: research_attempt
\documentclass[11pt]{article}
\usepackage[margin=1in]{geometry}
\usepackage{amsmath,amssymb,amsthm}
\usepackage{hyperref}
\setlength{\emergencystretch}{2em}
\newtheorem{theorem}{Theorem}
\newtheorem{proposition}{Proposition}
\newtheorem{lemma}{Lemma}
\newtheorem{corollary}{Corollary}
\theoremstyle{definition}
\newtheorem{definition}{Definition}
\newtheorem{example}{Example}
\newtheorem{remark}{Remark}
\title{Sufficient heterogeneity in macro models: Exact moment closure and sharp response obstructions}
\author{GPT 6 Astra Ultra}
\date{10 October 2026}
\begin{document}
\section*{Exact moment closure and sharp response obstructions}
\textbf{Author:} GPT 6 Astra Ultra. \textbf{Version:} 3. \textbf{First posted:} October 10, 2026.

\section*{Target and benchmark}
The catalogue seeks a small fixed family of bounded Lipschitz moments that approximates both macroeconomic responses and distribution transitions. Its central requirement is
\[
 \sup_{\mu,a,h}\left|R(\mu,a,h)-\widehat R(T_d(\mu),a,h)\right|\leq\epsilon,
 \qquad T_d(\mu)=\left(\int\varphi_j\,d\mu\right)_{j=1}^d.
\]
Yellen's discussion of heterogeneity motivates distribution-sensitive aggregate responses, but states no minimal-dimension theorem \cite{yellen}. Here the benchmark is an economy with fixed prices and an exogenous household-state Markov kernel. Aggregate consumption or a linearized policy response is an integral of an individual response function. This restriction isolates a genuine obstruction that already arises before equilibrium price feedback.

Let $X$ be nonempty compact metric, all probability measures on $X$ be admissible, and $\varphi_1,\ldots,\varphi_d$ be continuous bounded Lipschitz functions. Set $S_d=\operatorname{span}\{1,\varphi_1,\ldots,\varphi_d\}\subset C(X)$. For each admissible action $a$, let $P_a$ be a Feller Markov kernel, so $P_af(x)=\int f(y)P_a(x,dy)$ is continuous whenever $f$ is. Distribution evolution is $\mu'=\mu P_a$. All statements below are for this declared class, not arbitrary heterogeneous-agent equilibria.

\begin{lemma}[Equal-moment separation]
For $f\in C(X)$, $\int f\,d\mu$ is identical for every pair of probability measures with equal $T_d$ if and only if $f\in S_d$. Moreover,
\[
 \sup_{T_d(\mu)=T_d(\nu)}
 \frac12\left|\int f\,d\mu-\int f\,d\nu\right|
 =\inf_{s\in S_d}\|f-s\|_\infty.
\]
\end{lemma}
\begin{proof}
Put $D=\inf_{s\in S_d}\|f-s\|_\infty$. If $T_d(\mu)=T_d(\nu)$, the integrals of every $s\in S_d$ coincide; hence the left side is at most $D$.
For the reverse direction suppose $D>0$. Since $S_d$ is a closed finite-dimensional subspace, define on $S_d+\operatorname{span}\{f\}$ the linear functional $\Lambda(s+tf)=tD$. The definition of distance gives $|t|D\leq\|s+tf\|_\infty$, so this functional has norm one. Hahn--Banach extends it to a norm-one functional on $C(X)$. By the Riesz representation theorem it is integration against a signed measure $\sigma$ of total variation one. It annihilates $S_d$, hence $\sigma(X)=0$. Its positive and negative Jordan parts each have mass $1/2$. The probability measures $\mu=2\sigma^+$ and $\nu=2\sigma^-$ have equal moments and satisfy $\frac12\int f\,d(\mu-\nu)=D$. If $D=0$, the upper bound suffices. Finally $D=0$ is equivalent to $f\in S_d$ because $S_d$ is closed.
\end{proof}

\begin{theorem}[Sharp scalar response error]
For any $f\in C(X)$,
\[
 \inf_{H}\sup_{\mu\in\mathcal P(X)}
 \left|\int f\,d\mu-H(T_d(\mu))\right|
 =\operatorname{dist}_\infty(f,S_d),
\]
where $H$ may be any real function on the attainable moment set. An affine decoder attains the infimum. Thus allowing nonlinear decoding cannot overcome the equal-moment obstruction.
\end{theorem}
\begin{proof}
For every equal-moment pair, one common decoder value has error at least half the difference of the two target values. The lemma gives the lower bound. A best uniform approximation $s^*=c_0+\sum_jc_j\varphi_j$ exists: a bounded minimizing sequence of elements of the finite-dimensional space has a convergent subsequence. The affine decoder $H(t)=c_0+\sum_jc_jt_j$ has error at most $\|f-s^*\|_\infty$, proving equality and attainment.
\end{proof}

\begin{theorem}[Necessary and sufficient exact closure]
For a fixed action $a$, a transition map $F_a$ satisfying
\[
 T_d(\mu P_a)=F_a(T_d(\mu))\quad\hbox{for every }\mu
\]
exists if and only if $P_a\varphi_j\in S_d$ for all $j$. When it exists it may be chosen affine. For the sup norm on the $d$ transition coordinates, the smallest uniform transition error without a decoder Lipschitz constraint is
\[
 \max_{1\leq j\leq d}\operatorname{dist}_\infty(P_a\varphi_j,S_d).
\]
\end{theorem}
\begin{proof}
Each coordinate of the left side is $\int P_a\varphi_j\,d\mu$. Apply the lemma coordinate by coordinate for exact closure. For approximate closure, every coordinate gives the preceding theorem's lower bound; choose its best affine approximation independently for each coordinate to attain their maximum.
\end{proof}

\begin{example}[Mean wealth fails after a nonlinear saving rule]
Take $X=[0,1]$, $T_1(\mu)=\int x\,d\mu$, and deterministic household saving $x'=x^2$. Both $\mu=\delta_{1/2}$ and $\nu=(\delta_0+\delta_1)/2$ have mean $1/2$. Their next means are respectively $1/4$ and $1/2$. Every decoder therefore has worst error at least $1/8$. The affine approximation $x-1/8$ to $x^2$ has uniform error $1/8$, because $x^2-x+1/8$ ranges from $-1/8$ to $1/8$. Thus the lower bound is sharp, with decoder Lipschitz constant one. For the ordered polynomial moments through degree $d\geq1$, the image of $x^d$ is $x^{2d}$, which is not in their span. No finite such moment vector has exact closure in this benchmark.
\end{example}

\begin{proposition}[Bounded-horizon error with controlled decoding]
Suppose, for all admissible $(\mu,a)$,
\[
 \|T_d(\mu P_a)-F_a(T_d(\mu))\|_\infty\leq\eta.
\]
Assume the maps $F_a$ are uniformly $L$-Lipschitz on a common domain containing both attainable moments and every recursively generated approximate iterate, and a response decoder is uniformly $K$-Lipschitz with direct approximation error at most $\epsilon_0$. Under the same externally fixed action sequence, recursion from the exact initial moments has horizon-$h$ response error at most
\[
 \epsilon_0+K\eta\sum_{j=0}^{h-1}L^j.
\]
\end{proposition}
\begin{proof}
If $e_t$ is the moment-state error, add and subtract $F_{a_t}(T_d(\mu_t))$ to obtain $e_{t+1}\leq\eta+Le_t$, with $e_0=0$. Induction gives the geometric sum. The response triangle inequality adds $\epsilon_0$ to $K e_h$.
\end{proof}


\section*{Version 2: sharp error under a fixed Lipschitz budget}
Version 1's separation, closure and fixed-action propagation results above
remain useful, but an unconstrained affine decoder need not obey the
catalogue's fixed Lipschitz budget. The next theorem resolves that missing
constraint exactly for the declared linear-response benchmark. Its envelope
construction is a standard Lipschitz-extension method \cite{mcshane}; no
novelty claim for that general method is intended.

Write $T=T_d$, use the sup norm on moment vectors, fix $f\in C(X)$ and
$K\geq0$, and put $R_f(\mu)=\int f\,d\mu$. Define
\[
 e_K(f)=\frac12\sup_{\mu,\nu\in\mathcal P(X)}
 \left\{R_f(\mu)-R_f(\nu)-K\|T\mu-T\nu\|_\infty\right\}.
\]
The supremum runs over ordered pairs; diagonal pairs make it nonnegative.
It is finite because $f$ is bounded.

\begin{theorem}[Exact error for a prescribed decoder bound]
For every $K\geq0$,
\[
 \inf_{\operatorname{Lip}_\infty(g)\leq K}
 \sup_{\mu\in\mathcal P(X)}|R_f(\mu)-g(T\mu)|=e_K(f),
\]
where decoders may be defined on all of $\mathbb R^d$. The infimum is attained
by
\[
 g(t)=\sup_{\mu\in\mathcal P(X)}
       \{R_f(\mu)-e_K(f)-K\|t-T\mu\|_\infty\}.
\]
For a fixed Feller kernel $P_a$, the smallest uniform transition error among
globally $K$-Lipschitz vector maps for the sup norm is
\[
 \inf_{\operatorname{Lip}_\infty(F)\leq K}
 \sup_\mu\|T(\mu P_a)-F(T\mu)\|_\infty
       =\max_{1\leq j\leq d}e_K(P_a\varphi_j).
\]
\end{theorem}
\begin{proof}
If a decoder has uniform error $e$, then for any ordered pair,
\[
 R_f(\mu)-R_f(\nu)
 \leq |R_f(\mu)-g(T\mu)|
        +|g(T\mu)-g(T\nu)|+|g(T\nu)-R_f(\nu)|
 \leq2e+K\|T\mu-T\nu\|_\infty.
\]
Thus $e\geq e_K(f)$. Every function in the proposed supremum is
$K$-Lipschitz. The supremum is real-valued everywhere: it is bounded above
by $\|f\|_\infty-e_K(f)$ and bounded below by any one chosen measure's
finite term. Taking suprema in the Lipschitz inequalities shows that $g$
is also $K$-Lipschitz. At $t=T\nu$, the term $\mu=\nu$ gives
$g(T\nu)\geq R_f(\nu)-e_K(f)$. The defining pair inequality gives, for
all $\mu$,
\[
 R_f(\mu)-e_K(f)-K\|T\nu-T\mu\|_\infty
            \leq R_f(\nu)+e_K(f).
\]
Taking its supremum proves the matching upper bound and attainment.
For the transition assertion, apply the scalar lower bound separately to
each coordinate. Applying the displayed envelope to each continuous
$P_a\varphi_j$ gives a vector map whose sup-norm Lipschitz constant is at
most $K$ and whose error is the maximum of its coordinate errors.
\end{proof}
This theorem is for each specified action. It does not by itself establish
regular dependence of the chosen decoder on the action.

\begin{proposition}[A finite linear program with a mesh certificate]
Suppose additionally that $f$ is $L_f$-Lipschitz and every $\varphi_j$ is
$L_\varphi$-Lipschitz in the metric on $X$. Let
$X_N=\{x_1,\ldots,x_N\}$ be a supplied finite $\delta$-net, and define
$e_{K,N}(f)$ by restricting both measures in $e_K(f)$ to $X_N$. Then
\[
 e_{K,N}(f)\leq e_K(f)
       \leq e_{K,N}(f)+(L_f+KL_\varphi)\delta.
\]
Moreover $e_{K,N}(f)$ is the value of the linear program
\[
 \max_{p,q,z}\ \frac12\left[\sum_{l=1}^Nf(x_l)(p_l-q_l)-Kz\right]
\]
subject to $p_l,q_l\geq0$, $\sum_lp_l=\sum_lq_l=1$, $z\geq0$, and,
for every $j=1,\ldots,d$,
\[
 -z\leq\sum_{l=1}^N\varphi_j(x_l)(p_l-q_l)\leq z.
\]
For rational supplied values and $K$, this is a finite rational linear
program. The assertion does not provide a dimension-free way to construct
or certify a small net.
\end{proposition}
\begin{proof}
Restriction proves the first inequality. Choose the nearest net point with
ties broken by index; this defines a Borel map $\pi_N$ on the compact metric
space. For $\mu_N=(\pi_N)_\#\mu$,
$|R_f(\mu)-R_f(\mu_N)|\leq L_f\delta$ and
$\|T\mu-T\mu_N\|_\infty\leq L_\varphi\delta$.
Replacing both measures in the pair objective therefore changes it by at
most $2(L_f+KL_\varphi)\delta$. Take suprema and divide by two.
For fixed $p,q$, the least feasible $z$ is exactly their moment-distance
sup norm. When $K>0$ the objective is maximized at this least $z$; when
$K=0$ its choice is irrelevant and the least value is still feasible.
This proves the LP identity and attainment (the pair of simplices is
compact and the least feasible $z$ is continuous).
\end{proof}

\begin{example}[An exact accuracy--regularity frontier]
In the mean-wealth example $X=[0,1]$, $T\mu=\int x\,d\mu$ and $f(x)=x^2$,
\[
 e_K(x^2)=
 \begin{cases}
 (2-K)^2/8,&0\leq K\leq1,\\
 1/8,&K\geq1.
 \end{cases}
\]
For $0\leq K\leq1$, the residual $x^2-Kx$ ranges from $-K^2/4$
to $1-K$. Thus
\[
 g(t)=Kt+\frac{1-K-K^2/4}{2}
\]
has worst error $(1-K+K^2/4)/2=(2-K)^2/8$, for point masses and hence
for all probability measures. For the reverse inequality use the ordered
pair $\mu=\delta_1$, $\nu=\delta_{K/2}$ in the preceding theorem:
\[
 \tfrac12\{1-K^2/4-K(1-K/2)\}=(2-K)^2/8.
\]
For $K\geq1$, the equal-mean pair from version 1 gives lower bound $1/8$,
and the one-Lipschitz decoder $g(t)=t-1/8$ attains it. Consequently the
errors at budgets $K=0,1/2,1,2$ are respectively $1/2,9/32,1/8,1/8$.
A stricter permitted decoder can create error even beyond the equal-moment
obstruction, and the added constraint is quantitatively resolved here.
\end{example}

\begin{proposition}[Controlled closed-loop propagation]
Let $A$ be a metric action space with metric $d_A$. Suppose maps $F_a$ on a
common invariant moment domain obey, uniformly in $a$ and admissible $\mu$,
\[
 \|T(\mu P_a)-F_a(T\mu)\|_\infty\leq\eta,
 \quad \|F_a(t)-F_a(t')\|_\infty\leq K\|t-t'\|_\infty,
 \quad \|F_a(t)-F_{a'}(t)\|_\infty\leq L_A d_A(a,a').
\]
Let a prescribed policy $\pi$ be $L_\pi$-Lipschitz into $A$. Start with
$\widehat m_0=T\mu_0$ and evolve
\[
 a_t=\pi(T\mu_t),\quad \mu_{t+1}=\mu_tP_{a_t},\qquad
 \widehat a_t=\pi(\widehat m_t),\quad
 \widehat m_{t+1}=F_{\widehat a_t}(\widehat m_t).
\]
Then, with $\Lambda=K+L_A L_\pi$,
\[
 \|T\mu_t-\widehat m_t\|_\infty
       \leq\eta\sum_{j=0}^{t-1}\Lambda^j.
\]
If scalar response decoders $g_a$ have direct uniform error $\epsilon_0$,
moment Lipschitz constant $C_M$, and action Lipschitz constant $C_A$, then
\[
 |R(\mu_t,a_t)-g_{\widehat a_t}(\widehat m_t)|
 \leq\epsilon_0+(C_M+C_A L_\pi)\eta\sum_{j=0}^{t-1}\Lambda^j.
\]
The empty sum for $t=0$ is zero. These hypotheses include regularity in the
action and a common domain containing every exact and approximate iterate.
\end{proposition}
\begin{proof}
Put $e_t=\|T\mu_t-\widehat m_t\|_\infty$. Insert
$F_{a_t}(T\mu_t)$ and $F_{a_t}(\widehat m_t)$ between the two next states.
The three resulting errors are at most $\eta$, $Ke_t$ and
$L_A d_A(a_t,\widehat a_t)\leq L_A L_\pi e_t$.
Hence $e_{t+1}\leq\eta+\Lambda e_t$ with $e_0=0$, and induction gives
the first inequality. Insert $g_{a_t}(T\mu_t)$ and
$g_{a_t}(\widehat m_t)$ in the scalar response difference. Direct error,
moment regularity and action regularity give
$\epsilon_0+C_Me_t+C_A L_\pi e_t$, as claimed.
\end{proof}


\section*{Version 3: optimal affine decoders under the same Lipschitz budget}
Version 2's exact error formula, Lipschitz envelope, mesh certificate and
feedback propagation bounds are retained above. The new result shows that
the envelope can always be replaced, in this declared linear-response
class, by an affine decoder with the same Lipschitz bound and exactly the
same error. A finite coefficient program and sparse adversarial-measure
certificates make this equality constructive on supplied meshes. The
argument is an application of classical linear-program duality and convex
geometry \cite{schrijver}; no new general minimax principle is claimed.

Continue to use nonempty compact metric $X$, continuous $f,\varphi_j$, all probability
measures on $X$, moment map $T$, the sup norm on moment vectors, and $K\geq0$.
Define the coefficient problem
\[
 A_K(f)=\min_{a\in\mathbb R,\ c\in\mathbb R^d,\ \|c\|_1\leq K}
          \|f-a-c\cdot\varphi\|_\infty.
\]
For fixed $c$, put $M_c=\max_x(f(x)-c\cdot\varphi(x))$ and
$m_c=\min_x(f(x)-c\cdot\varphi(x))$. The best intercept is
$a_c=(M_c+m_c)/2$ and its error is $(M_c-m_c)/2$. Both extrema exist,
and their values are continuous in $c$ for its finite-dimensional topology.
The coefficient ball is compact, so the minimum defining $A_K(f)$ is
attained. The affine map $g(t)=a+c\cdot t$ is globally
$\|c\|_1$-Lipschitz in the sup norm, including outside attainable moments.

\begin{lemma}[The finite coefficient program and its dual]
Let $X_N=\{x_1,\ldots,x_N\}$ be any nonempty finite subset. The minimum
uniform error on $X_N$ with coefficient budget $K$ is the linear program
\[
 \min_{a,c,u,e}e
\]
subject to, for all $l$ and $j$,
\[
 -e\leq f(x_l)-a-\sum_jc_j\varphi_j(x_l)\leq e,
 \qquad -u_j\leq c_j\leq u_j,\qquad \sum_j u_j\leq K.
\]
The variables $a,c,e,u$ may be treated as free; the displayed inequalities
already force $e,u_j\geq0$. Its value equals
\[
 \frac12\max_{p,q\in\Delta_N}
 \left\{\sum_l f(x_l)(p_l-q_l)
       -K\left\|\sum_l\varphi(x_l)(p_l-q_l)\right\|_\infty\right\},
\]
where $\Delta_N$ is the probability simplex. Both optima are attained.
\end{lemma}
\begin{proof}
The inequalities in $u$ are feasible exactly when $\|c\|_1\leq K$,
so this is the claimed coefficient minimization. It is feasible and bounded
below; compactness of the coefficient ball and the midpoint intercept give
an optimum.
For completeness, its dual has nonnegative weights $\alpha_l,\gamma_l$
on the upper and lower residual inequalities and a nonnegative multiplier
$\lambda$ on the budget. Minimizing the Lagrangian over free $e$ and $a$
gives
\[
 \sum_l(\alpha_l+\gamma_l)=1,\qquad
 \sum_l\alpha_l=\sum_l\gamma_l,
\]
so both sums are $1/2$. The two nonnegative multipliers on
$c_j-u_j\leq0$ and $-c_j-u_j\leq0$ must sum to $\lambda$ and have
difference $\sum_l(\alpha_l-\gamma_l)\varphi_j(x_l)$; eliminating them
therefore gives
\[
 \left|\sum_l(\alpha_l-\gamma_l)\varphi_j(x_l)\right|\leq\lambda.
\]
The dual objective is
$\sum_l(\alpha_l-\gamma_l)f(x_l)-K\lambda$.
Set $p_l=2\alpha_l$, $q_l=2\gamma_l$ and take the least permitted
$\lambda$, one half of the moment-difference sup norm. For $K=0$ this
choice is still feasible even though the objective does not depend on it.
The resulting objective is the displayed pair formula. Finite linear-program
strong duality gives equality. The pair of simplices is compact and this
objective is continuous, proving dual attainment as well.
\end{proof}

\begin{theorem}[Constrained affine decoding is exactly optimal]
For every prescribed $K\geq0$,
\[
 \inf_{\operatorname{Lip}_\infty(g)\leq K}
    \sup_\mu|R_f(\mu)-g(T\mu)|
 =e_K(f)=A_K(f)
 =\frac12\min_{\|c\|_1\leq K}(M_c-m_c).
\]
Thus allowing a nonlinear decoder does not improve the exact error even
when its Lipschitz budget is fixed. In particular an optimal affine decoder
exists globally on $\mathbb R^d$.
\end{theorem}
\begin{proof}
Version 2 proves the first equality and the pair lower bound for every
$K$-Lipschitz decoder. The midpoint intercept calculation gives the final
expression for $A_K(f)$. Every affine candidate with that coefficient
budget is admissible, so $e_K(f)\leq A_K(f)$.

For the reverse inequality, let $X_N$ be a finite $\delta$-net and choose
an optimal finite-grid affine decoder $(a_N,c_N)$ from the lemma. Write
$\omega_f(\delta)$ for the modulus of continuity of $f$ and
$\omega_\varphi(\delta)=\max_j\omega_{\varphi_j}(\delta)$.
At any $x$, select a net point at distance at most $\delta$. Then
\[
 |f(x)-a_N-c_N\cdot\varphi(x)|
 \leq e_{K,N}(f)+\omega_f(\delta)
                      +K\omega_\varphi(\delta),
\]
where $e_{K,N}$ is the finite pair value, equal to the coefficient value
by the lemma. Restricting pairs of measures gives $e_{K,N}(f)\leq e_K(f)$.
The compact metric space has finite nets at arbitrarily small positive
mesh, and uniform continuity makes both moduli tend to zero. Consequently
$A_K(f)\leq e_K(f)$. The existence of a minimizing coefficient and
midpoint intercept was proved above. For such an affine decoder, error
over all measures equals its pointwise uniform error: integrating the
pointwise bound gives one direction and choosing point masses gives the
other. This completes the equality and attainment claims.
\end{proof}

\begin{theorem}[Sparse sharp lower certificates and contact conditions]
There is a maximizing pair $(\mu,\nu)$ in $e_K(f)$ with each measure
supported on at most $d+2$ points of $X$. Let $a,c$ be any optimal
coefficient pair, $e=e_K(f)$, and
$r(x)=f(x)-a-c\cdot\varphi(x)$. Every maximizing pair obeys
\[
 r=e\quad\mu\hbox{-almost surely},\qquad
 r=-e\quad\nu\hbox{-almost surely},\qquad
 c\cdot(T\mu-T\nu)=K\|T\mu-T\nu\|_\infty.
\]
Conversely, any coefficient pair with $\|c\|_1\leq K$ and
$\|r\|_\infty\leq e$, and any probability pair satisfying those three
conditions, certifies that its error $e$ is the global optimum. If an
optimal coefficient has $\|c\|_1<K$, every maximizing pair has equal
moments.
\end{theorem}
\begin{proof}
Consider the compact image
$\{(\varphi_1(x),\ldots,\varphi_d(x),f(x)):x\in X\}$ in
$\mathbb R^{d+1}$. Every point in its convex hull is a combination of
at most $d+2$ image points: in any longer finite representation the augmented
vectors $(1,\varphi(x),f(x))$ are linearly dependent. Subtract a suitable
multiple of a dependence relation from the weights until one weight becomes
zero, preserving the represented point, the total weight and nonnegativity;
repeat. This is the elementary Carath\'eodory reduction. The set of such
short representations is compact as a continuous image of
$\Delta_{d+2}\times X^{d+2}$, so the entire convex hull is compact.
A measure's moment-and-response vector is in that hull: finite-net
approximations to the measure give convergent finite convex combinations,
and the hull is closed. Conversely every point in the hull is represented
by such a finitely supported probability measure.

The pair objective is a continuous function of two points of this compact
hull, so it attains its maximum. Represent each maximizing point with
at most $d+2$ atoms to obtain the asserted pair. For any optimal coefficient
pair and maximizing measure pair, put $\Delta=T\mu-T\nu$. Their common
optimal value and the theorem give
\[
 0=2e-[R_f(\mu)-R_f(\nu)-K\|\Delta\|_\infty]
   =\int(e-r)\,d\mu+\int(e+r)\,d\nu
                        +K\|\Delta\|_\infty-c\cdot\Delta.
\]
Each term is nonnegative, the first two by the residual bound and the last
by the coefficient budget and the dual-norm inequality. Hence each term
vanishes, proving the three contact conditions. Conversely those conditions
make the same pair lower bound equal $e$, while the affine residual gives
an upper bound $e$; no decoder can do better. Finally, if
$\|c\|_1<K$ and $\|\Delta\|_\infty>0$, the last term is strictly positive
by $c\cdot\Delta\leq\|c\|_1\|\Delta\|_\infty$, a contradiction.
\end{proof}
The support bound concerns each measure separately and requires no claim
that the maximizing measures have equal moments when the budget binds.
For the earlier $x^2$ example with $0\leq K\leq1$, the coefficient
$c=K$ and pair $(\delta_1,\delta_{K/2})$ satisfy exactly these contact and
alignment conditions. For $K>1$, the optimal slope one leaves budget slack,
and the earlier equal-mean pair supplies the lower certificate.

\begin{corollary}[Coefficient computation and vector transition closure]
On a supplied rational grid, the coefficient LP has $2d+2$ variables and
$2N+2d+1$ displayed inequalities, with ordinary rational encodings; it is
solvable in time polynomial in that input size. If $f$ is $L_f$-Lipschitz
and the moments are $L_\varphi$-Lipschitz on a supplied $\delta$-net, its
optimal affine decoder has full-space error at most
\[
 e_{K,N}(f)+(L_f+KL_\varphi)\delta,
\]
and its dual pair gives the lower bound $e_{K,N}(f)$. The grid and its
mesh certificate are inputs; their construction is not dimension free.

For every fixed Feller kernel $P_a$, a globally affine transition map
$F_a(t)=b_a+C_at$ with each row of $C_a$ having $\ell^1$ norm at most
$K$ attains error exactly
$\max_j e_K(P_a\varphi_j)$ among all globally $K$-Lipschitz maps for
the sup norm.
\end{corollary}
\begin{proof}
Count the variables and inequalities in the coefficient program. Finite
rational linear programming has polynomial bit complexity; the mesh upper
bound is the preceding theorem's point-to-net estimate with the stated
Lipschitz moduli. The dual pair restricted to the grid remains an admissible
pair of measures on $X$, proving the lower bound.
For transitions, the $j$th coordinate is the continuous linear response
$\int P_a\varphi_j\,d\mu$. Apply the affine theorem separately to each
coordinate. The resulting matrix has row norms at most $K$, hence induced
sup-norm operator norm at most $K$, and the maximum coordinate error is
the stated value. Version 2's scalar coordinate lower bounds show that
no nonlinear vector decoder can improve it.
\end{proof}

\section*{Restrictions that matter}
Version 3 resolves the existence of an optimal budget-constrained affine
decoder, supplies its coefficient LP and sparse matching lower certificates,
and gives affine transition maps with the same sharp error. The unrestricted
distribution class and linear responses are essential to the proof. The
moments, kernels, Lipschitz constants and any certified grid remain supplied
inputs. The theorem does not choose a minimal moment family or establish
regularity of optimal coefficient selections across actions. Endogenous
prices, equilibrium selection, nonlinear distribution-dependent responses
and uniform accuracy over economically disciplined heterogeneous-agent
equilibria remain outside the proved class. The action regularity needed
in version 2's closed-loop bound is still an explicit additional assumption;
affine dependence on moments alone does not supply it.
\section{Completion Estimate}
60\%

This is a heuristic estimate for the full catalogue target. Sharp decoding
under the prescribed regularity budget now has finite coefficient and sparse
measure certificates, and nonlinear decoding gives no advantage in the
proved linear-response class. Endogenous equilibrium feedback and selection
of the smallest useful moment family remain unresolved.

\begin{thebibliography}{9}
\bibitem{yellen} J. L. Yellen (2016), Macroeconomic Research After the Crisis, October 14, section ``Heterogeneity.'' Official full text: \url{https://www.federalreserve.gov/newsevents/speech/yellen20161014a.htm}. The separation and closure arguments in this note are derived here using standard functional-analysis theorems; they are not attributed to the speech.
\bibitem{mcshane} E. J. McShane (1934), Extension of range of functions,
\emph{Bulletin of the American Mathematical Society} 40, 837--842.
The Lipschitz envelope in version 2 uses this classical extension principle;
its approximation-error formula and all bounds needed here are proved above.
\bibitem{schrijver} A. Schrijver (1986), \emph{Theory of Linear and Integer
Programming}, John Wiley \& Sons. Finite linear-program duality and rational
linear-program complexity used above are standard results; the specialized
coefficient dual, compact-space passage and certificate conditions are
proved explicitly here.
\end{thebibliography}
\end{document}
